Un ensemble de droites du plan est dit en \emph{position générale} s'il n'y en a pas deux parallèles ni trois
qui passent par un même point. Un ensemble de droites en position générale découpe le plan en régions,
certaines ayant une aire finie. Ces dernières sont appelées les \emph{régions finies de l'ensemble}. Prouver
que pour tout $n$ suffisamment grand, dans tout ensemble de $n$ droites en position générale, il est possible
de colorer en bleu au moins $\sqrt{n}$ de ses droites de sorte qu'il n'y ait aucune région finie de cet
ensemble de bord entièrement bleu.

\emph{Note :} Des résultats pour lesquels $\sqrt{n}$ serait remplacé par $c\sqrt{n}$ seront valorisés selon la
valeur de la constante $c$.
